\chapter{The Matrix-Free Kernel in Detail}
\label{app:kernel}

\chapterorient{\Cref{ch:matrixfree} factored the finite-element operator into five
element-local stages, $\mat{A}\vv = G^{\mathsf T}B_{\diffop}^{\mathsf T}\,D\,
B_{\diffop}\,G\,\vv$, and promised this appendix as the deep version. Here we
deliver it. We take each of the five stages apart down to its exact tensor shape,
write out the index maps that $G$ and $G^{\mathsf T}$ actually perform, show how
the differential operator $\diffop$ selects the basis mode and shapes the
component axis $C$, and build the geometry factor that the diagonal $D$ consumes
from the Jacobian metric up. We separate, in full, the build-time stratum (the
geometry computed once per mesh) from the run-time stratum (the action recomputed
on every apply); we count the operations sum-factorization saves and present it as
a contraction reordering; and we draw the libCEED library boundary precisely,
distinguishing the kernel \emph{API} we call from the kernel \emph{implementation}
---the contraction chain of this appendix---that realizes it. By the end the whole
operator stands exposed as a composition of tensor contractions over named axes:
the GPU-native form, with no sparse matrix anywhere, that the book's
implementation view (\cref{ch:synthlib}) consumes. This is where the mathematics
meets the silicon.}

\section{Scope: the chain this appendix opens}
\label{sec:appk-scope}

\Cref{ch:matrixfree} established the central factorization and verified it against
an assembled matrix on a two-element 1-D mesh. We restate it once, because every
section below is a magnification of one of its factors:
\begin{equation}
  \mat{A}\,\vv
   \;=\;
   \underbrace{G^{\mathsf T}}_{\text{scatter}}\;
   \underbrace{B_{\diffop}^{\mathsf T}}_{\text{integrate}}\;
   \underbrace{D}_{\text{weight}}\;
   \underbrace{B_{\diffop}}_{\text{evaluate}}\;
   \underbrace{G}_{\text{gather}}\;
   \vv .
  \label{eq:appk-chain}
\end{equation}
The five factors carry a tensor through a sequence of named shapes,
\begin{equation}
\begin{aligned}
  \texttt{[(N: \dots)]}
  &\xrightarrow{\,G\,}
  \texttt{[(E, L)]}
  \xrightarrow{\,B_{\diffop}\,}
  \texttt{[(E, P, C)]}
  \xrightarrow{\,D\,}
  \texttt{[(E, P, C')]} \\
  &\xrightarrow{\,B_{\diffop}^{\mathsf T}\,}
  \texttt{[(E, L)]}
  \xrightarrow{\,G^{\mathsf T}\,}
  \texttt{[(N: \dots)]},
\end{aligned}
  \label{eq:appk-shapes}
\end{equation}
where the five axes are the element-local family fixed in \cref{ch:refelement} and
catalogued in \cref{app:shapes}: $N$ the global true-degree-of-freedom count, $E$
the element count, $L$ the local dofs per element, $P$ the quadrature points per
element, $C$ the value/derivative components carried at a quadrature point, and
$G$ the geometry-metric components stored per point. \Cref{tab:appk-axes} is the
working legend for this appendix; \cref{app:shapes} carries the full algebra.

\begin{table}[t]
\centering
\caption[The five element-local axes, with extents]{The element-local tensor axes
used throughout this appendix, with their typical extents on a degree-$p$,
$d$-dimensional tensor-product element. The stratum column---built once per
$(\text{mesh},\text{order})$ versus produced afresh on every apply---is the subject
of \cref{sec:appk-strata}. The global axis $N$ is \emph{not} an element-local axis;
it is the flat solver vector, and $G$ is the boundary between it and the
$(E,L)$ vocabulary.}
\label{tab:appk-axes}
\small
\setlength{\tabcolsep}{4.5pt}
\renewcommand{\arraystretch}{1.25}
\begin{tabular}{@{}cllcl@{}}
\toprule
axis & name & meaning & extent & stratum \\
\midrule
$N$ & global dofs & true dofs of the global vector & --- & solver vector \\
$E$ & element count & mesh elements assembled over & --- & build-time \\
$L$ & local dofs & dofs local to one element & $(p{+}1)^d$ & build-time \\
$P$ & quad points & quadrature points per element & $(p{+}1)^d$ & build-time \\
$C$ & components & value/deriv.\ components per point & $1$ or $d$ & run-time \\
$G$ & geom comps & stored metric comps per point & $2+d^2$ & build-time \\
\bottomrule
\end{tabular}
\end{table}

\Cref{fig:appk-pipeline} is the deep version of \cref{ch:matrixfree}'s hero
diagram: the same five stages, but with the \emph{exact} shape on every wire,
including the build-time inputs each stage draws from below. The sections that
follow walk it stage by stage.

\begin{figure}[p]
\centering
\resizebox{\linewidth}{!}{%
\begin{tikzpicture}[node distance=7mm and 8mm]
  % input
  \node[data, text width=20mm] (in) {global vector\\\code{[(N: \dots)]}};
  \node[stage, right=of in, text width=19mm] (G)  {$G$\\\footnotesize gather\\element restrict};
  \node[stage, right=of G,  text width=19mm] (B)  {$B_{\diffop}$\\\footnotesize evaluate\\basis apply};
  \node[extern, right=of B, text width=19mm] (D)  {$D$\\\footnotesize weight\\\code{\#extern}};
  \node[stage, right=of D,  text width=19mm] (Bt) {$B_{\diffop}^{\mathsf T}$\\\footnotesize integrate};
  \node[stage, right=of Bt, text width=19mm] (Gt) {$G^{\mathsf T}$\\\footnotesize scatter-add};
  \node[data, right=of Gt, text width=20mm] (out) {global vector\\\code{[(N: \dots)]}};

  \draw[flow] (in) -- (G);  \draw[flow] (G) -- (B); \draw[flow] (B) -- (D);
  \draw[flow] (D) -- (Bt);  \draw[flow] (Bt) -- (Gt); \draw[flow] (Gt) -- (out);

  % exact shapes on every wire
  \node[font=\scriptsize\ttfamily, simGray, below=2.4mm of G.south]  {[(E, L)]};
  \node[font=\scriptsize\ttfamily, simGray, below=2.4mm of B.south]  {[(E, P, C)]};
  \node[font=\scriptsize\ttfamily, simGray, below=2.4mm of D.south]  {[(E, P, C')]};
  \node[font=\scriptsize\ttfamily, simGray, below=2.4mm of Bt.south] {[(E, L)]};
  \node[font=\scriptsize\ttfamily, simBlue, below=3mm of in.south, align=center] {action\\input};
  \node[font=\scriptsize\ttfamily, simBlue, below=3mm of out.south, align=center] {action\\output};

  % build-time inputs fed from below
  \node[data, fill=simBgGold, draw=simGold, below=15mm of G, text width=22mm]
    (restr) {\footnotesize\itshape restriction\\index map (build)};
  \node[data, fill=simBgGold, draw=simGold, below=15mm of B, text width=22mm]
    (basis) {\footnotesize\itshape tabulated basis\\$\hat B,\,\hat G$ (build)};
  \node[data, fill=simBgGold, draw=simGold, below=15mm of D, text width=24mm]
    (geom) {\footnotesize\itshape geom\_data \code{[(E,P,G)]}\\(build)};
  \draw[flow, simGold] (restr) -- (G);
  \draw[flow, simGold] (basis) -- (B);
  \draw[flow, simGold] (geom)  -- (D);
  % Bt and Gt reuse the same build-time tables (dotted back-refs)
  \draw[refedge, simGold] (basis.north) .. controls +(0,0.9) and +(0,-1.0) .. (Bt.south);
  \draw[refedge, simGold] (restr.north) .. controls +(0,1.1) and +(0,-1.2) .. (Gt.south);

  % brackets: element-local middle
  \draw[decorate, decoration={brace, amplitude=4pt}, simRust]
    ([yshift=7mm]G.north west) -- ([yshift=7mm]Gt.north east)
    node[midway, above=4pt, font=\scriptsize, simRust]
    {element-local: block-diagonal in $E$, fully parallel};
\end{tikzpicture}%
}
\caption[The five-stage pipeline with exact tensor shapes]{The matrix-free
operator of \cref{eq:appk-chain}, drawn with the \emph{exact} element-local shape
on every wire (the deep version of \cref{ch:matrixfree}'s pipeline figure). A flat
global vector \code{[(N: \dots)]} enters; the gather $G$ produces per-element local
blocks \code{[(E, L)]}; basis evaluation $B_{\diffop}$ lifts them to
\code{[(E, P, C)]} values at the quadrature points; the center stage $D$, the
\code{\#extern} libCEED boundary (violet, \cref{sec:appk-extern}), weights each
point against the precomputed \code{geom\_data}, possibly changing the component
count to $C'$; $B_{\diffop}^{\mathsf T}$ integrates back to \code{[(E, L)]}; and
$G^{\mathsf T}$ scatters into a fresh global vector. The three build-time tables
(gold, fed from below)---the restriction index map, the tabulated basis $\hat B$
and reference-gradient $\hat G$, and the geometry factor---are computed once and
reused by both the forward stages and their transposes (dotted back-references).
The element-local middle (rust brace) is block-diagonal in $E$.}
\label{fig:appk-pipeline}
\end{figure}

\section{The element-local tensor, axis by axis}
\label{sec:appk-anatomy}

Before opening the stages we fix the three tensors they move between, because
their axes are the entire vocabulary. The flat solver vector is rank-1,
\code{[(N: \dots)]}; the moment the gather runs we acquire an explicit element
axis $E$ and the per-element structure under it. \Cref{fig:appk-anatomy} is the
anatomy chart.

\begin{figure}[t]
\centering
\begin{tikzpicture}[scale=1.0]
  % --- [E, L] ---
  \begin{scope}[shift={(-4.6,0)}]
    \foreach \k in {0,1,2}{
      \begin{scope}[shift={(0.12*\k,0.09*\k)}]
        \fill[simBgBlue, draw=simBlue] (-0.7,-0.5) rectangle (0.7,0.5);
        \foreach \y in {-0.25,0,0.25}{\draw[simBlue!40] (-0.7,\y)--(0.7,\y);}
      \end{scope}}
    \node[font=\scriptsize, simInk] at (-1.05,0.0) {$L$};
    \node[font=\scriptsize, simRust] at (0.55,0.78) {$E$};
    \node[font=\scriptsize\ttfamily, simInk] at (0.1,-1.0) {[(E, L)]};
    \node[font=\scriptsize\itshape, simGray, text width=26mm, align=center]
      at (0.1,-1.52) {per-element local dofs\\(restricted field)};
  \end{scope}
  \node[font=\large] at (-2.7,0) {$\xrightarrow{\,B_{\diffop}\,}$};
  % --- [E, P, C] ---
  \begin{scope}[shift={(-0.6,0)}]
    \foreach \k in {0,1,2}{
      \begin{scope}[shift={(0.12*\k,0.09*\k)}]
        \fill[simBgGreen, draw=simGreen] (-0.78,-0.5) rectangle (0.78,0.5);
        \foreach \y in {-0.17,0.17}{\draw[simGreen!40] (-0.78,\y)--(0.78,\y);}
        \foreach \x in {-0.26,0.26}{\draw[simGreen!40] (\x,-0.5)--(\x,0.5);}
      \end{scope}}
    \node[font=\scriptsize, simInk] at (-1.13,0.0) {$P$};
    \node[font=\scriptsize, simInk] at (0.0,-0.83) {$C$};
    \node[font=\scriptsize, simRust] at (0.62,0.78) {$E$};
    \node[font=\scriptsize\ttfamily, simInk] at (0.05,-1.0) {[(E, P, C)]};
    \node[font=\scriptsize\itshape, simGray, text width=28mm, align=center]
      at (0.05,-1.52) {$C$ values at\\$P$ quad points};
  \end{scope}
  \node[font=\large] at (1.55,0) {$\odot$};
  % --- [E, P, G]  geom_data ---
  \begin{scope}[shift={(3.7,0)}]
    \foreach \k in {0,1,2}{
      \begin{scope}[shift={(0.12*\k,0.09*\k)}]
        \fill[simBgGold, draw=simGold] (-0.78,-0.5) rectangle (0.78,0.5);
        \foreach \y in {-0.17,0.17}{\draw[simGold!50] (-0.78,\y)--(0.78,\y);}
        \foreach \x in {-0.26,0.26}{\draw[simGold!50] (\x,-0.5)--(\x,0.5);}
      \end{scope}}
    \node[font=\scriptsize, simInk] at (-1.13,0.0) {$P$};
    \node[font=\scriptsize, simInk] at (0.0,-0.83) {$G$};
    \node[font=\scriptsize, simRust] at (0.62,0.78) {$E$};
    \node[font=\scriptsize\ttfamily, simInk] at (0.05,-1.0) {[(E, P, G)]};
    \node[font=\scriptsize\itshape, simGray, text width=28mm, align=center]
      at (0.05,-1.52) {geom\_data\\(build-time)};
  \end{scope}
\end{tikzpicture}
\caption[Anatomy of the element-local tensor family]{The three element-local
tensors and the axes they share. Gather produces \code{[(E, L)]} (blue): an
element axis $E$ over per-element local-dof blocks of length $L$. Basis evaluation
$B_{\diffop}$ maps it to \code{[(E, P, C)]} (green): $C$ components at each of $P$
quadrature points, per element. The weight $D$ multiplies it pointwise ($\odot$)
against the precomputed geometry carrier \code{[(E, P, G)]} (gold). The element
axis $E$ (rust) threads through all three: every stage between gather and scatter
is block-diagonal in $E$. The $G$ axis is purely build-time storage; the $C$ axis
is the only run-time-shaped axis, and even it is fixed once the differential
operator $\diffop$ is chosen.}
\label{fig:appk-anatomy}
\end{figure}

\begin{notationbox}
The libCEED library, which owns the device kernel (\cref{sec:appk-extern}), names
two of these axes \emph{oppositely} to this book. Its basis constructor's
parameter \code{P} counts \emph{nodes per element}---our $L$---and its \code{Q}
counts \emph{quadrature points}---our $P$. This appendix uses $L$ for local dofs
and $P$ for quadrature points throughout; if you read the library source expect the
letters to swap. The binding is exact; only the labels differ.
\end{notationbox}

The $G$ axis deserves an early word because it is the one whose extent we have not
yet pinned. The geometry carrier stores $G = 2 + d^2$ components per quadrature
point: in three dimensions, $2 + 9 = 11$. The two scalar slots hold the
quadrature-weighted volume factor and an attribute (the material-region tag); the
$d^2$ slots hold the full metric matrix $J^{-1}J^{-\mathsf T}|\det J|$ that a
grad--grad term needs (\cref{sec:appk-geom}). A pure mass term uses only the volume
scalar, but the storage size is fixed by the largest term, so $G=2+d^2$ is the
allocation. This is a positive fact read off the library's build pass, not an
estimate: the assembler verifies \code{geom\_data\_size == 2 + space\_dim*dim}
before it runs.

\section{Stage 1: \texorpdfstring{$G$}{G}, element restriction (the gather)}
\label{sec:appk-gather}

The gather $G$ is the boundary between two vocabularies. On its left sits the
flat global vector \code{[(N: \dots)]}, a rank-1 list the solver owns; on its right
sit the per-element blocks \code{[(E, L)]} the rest of the pipeline runs on. $G$
is the map that crosses it.

\subsection{The index map}

$G$ performs no arithmetic. It is a pure \emph{indexed copy}, defined by a fixed
table---the \emph{element-to-dof map}---that records, for each element $e$ and each
of its $L$ local dof slots $\ell$, which global dof index $n = \iota(e,\ell)$ it
reads from. The action is simply
\begin{equation}
  (G\vv)[e,\ell] \;=\; \vv\bigl[\iota(e,\ell)\bigr],
  \qquad e\in\{0,\dots,E-1\},\ \ell\in\{0,\dots,L-1\}.
  \label{eq:appk-gather}
\end{equation}
For each element, $G$ walks that element's $L$ slots and copies in the global
values they point at. A global dof that several elements share---a vertex,
edge, or face dof on an inter-element boundary---is read by \emph{every} element
that touches it, so its value is \emph{copied out} into several local blocks. That
duplication is deliberate; it is what lets the middle stages run element by element
in ignorance of the sharing, and it is undone, by summation, at the far end of the
pipeline. \Cref{fig:appk-gatherscatter} draws the copy-out and the matching
sum-back.

\begin{figure}[t]
\centering
\begin{tikzpicture}[scale=1.0, every node/.style={font=\scriptsize}]
  % ---- gather (top): one shared global dof fans out ----
  \begin{scope}
    \node[font=\footnotesize\bfseries, simBlue] at (-2.4,1.1) {$G$ (gather)};
    % global slots
    \foreach \i/\lab in {0/$v_0$,1/$v_1$,2/$v_2$,3/$v_3$}{
      \node[draw=simGray, fill=simBgGray, minimum size=5.5mm] (g\i) at (\i*0.9,1.1) {\lab};}
    \node[simGray] at (1.35,1.7) {global \code{[(N)]}};
    % element blocks e0=(v0,v1) e1=(v1,v2) e2=(v2,v3)
    \foreach \k/\a/\b/\la/\lb in {0/0/0.9/$v_0$/$v_1$, 1/2.2/3.1/$v_1$/$v_2$, 2/4.4/5.3/$v_2$/$v_3$}{
      \node[draw=simBlue, fill=simBgBlue, minimum size=5.0mm] (e\k a) at (\a,-0.2) {\la};
      \node[draw=simBlue, fill=simBgBlue, minimum size=5.0mm] (e\k b) at (\b,-0.2) {\lb};
      \node[simRust] at ({(\a+\b)/2},-0.75) {$e_\k$};}
    % copy arrows for the shared dof v1 and v2 (fan-out)
    \draw[flow, simBlue] (g0.south) -- (e0a.north);
    \draw[flow, simBlue] (g1.south) -- (e0b.north);
    \draw[flow, simRust, thick] (g1.south) -- (e1a.north);   % v1 fans to e1 too
    \draw[flow, simBlue] (g2.south) -- (e1b.north);
    \draw[flow, simRust, thick] (g2.south) -- (e2a.north);   % v2 fans to e2 too
    \draw[flow, simBlue] (g3.south) -- (e2b.north);
    \node[simRust, font=\scriptsize\itshape] at (6.6,0.4) {shared dof};
    \node[simRust, font=\scriptsize\itshape] at (6.6,0.05) {copied out};
  \end{scope}
  % ---- scatter (bottom): contributions sum back ----
  \begin{scope}[shift={(0,-3.4)}]
    \node[font=\footnotesize\bfseries, simGreen] at (-2.4,-0.2) {$G^{\mathsf T}$ (scatter-add)};
    \foreach \k/\a/\b/\la/\lb in {0/0/0.9/$c^0_0$/$c^0_1$, 1/2.2/3.1/$c^1_0$/$c^1_1$, 2/4.4/5.3/$c^2_0$/$c^2_1$}{
      \node[draw=simBlue, fill=simBgBlue, minimum size=5.0mm] (s\k a) at (\a,0.5) {\la};
      \node[draw=simBlue, fill=simBgBlue, minimum size=5.0mm] (s\k b) at (\b,0.5) {\lb};}
    \foreach \i/\lab in {0/{$c^0_0$},1/{$c^0_1{+}c^1_0$},2/{$c^1_1{+}c^2_0$},3/{$c^2_1$}}{
      \node[draw=simGreen, fill=simBgGreen, minimum width=9mm, minimum height=5.5mm] (o\i) at (\i*1.55+0.0,-0.9) {\lab};}
    \draw[flow, simGreen] (s0a.south) -- (o0.north);
    \draw[flow, simGreen] (s0b.south) -- (o1.north);
    \draw[flow, simGreen] (s1a.south) -- (o1.north);
    \draw[flow, simGreen] (s1b.south) -- (o2.north);
    \draw[flow, simGreen] (s2a.south) -- (o2.north);
    \draw[flow, simGreen] (s2b.south) -- (o3.north);
    \node[simGreen, font=\scriptsize\itshape] at (6.6,-0.5) {summed};
    \node[simGreen, font=\scriptsize\itshape] at (6.6,-0.85) {at shared dofs};
  \end{scope}
\end{tikzpicture}
\caption[Gather copies out; scatter sums back]{The $G$/$G^{\mathsf T}$ transpose
pair on a 1-D mesh of three elements sharing two interior nodes. \emph{Top:} the
gather $G$ copies each global value into every element that references it; the
shared dofs $v_1,v_2$ (rust) fan out to two elements each. No arithmetic---pure
copying. \emph{Bottom:} the scatter-add $G^{\mathsf T}$ runs the table in reverse,
\emph{summing} the per-element contributions $c^e_\ell$ back into each global slot;
at a shared node the two elements' contributions add (e.g.\ $c^0_1+c^1_0$). The
copy-out/sum-back pair is the entire content of finite-element assembly, here
applied to a vector. Because $G^{\mathsf T}$ is the literal transpose of $G$,
$\inner{G^{\mathsf T}\vw}{\vv}=\inner{\vw}{G\vv}$ holds exactly---which is why the
assembled value is a \emph{sum} and not an average.}
\label{fig:appk-gatherscatter}
\end{figure}

\subsection{$G^{\mathsf T}$ is the scatter-add; $G^{\mathsf T}G$ is the multiplicity}

The transpose $G^{\mathsf T}$ runs the same index table backwards. Where $G$ read
$n=\iota(e,\ell)$ and copied $\vv[n]$ into slot $(e,\ell)$, the transpose takes a
per-element tensor $\vc\in\code{[(E, L)]}$ and \emph{adds} each local contribution
into the global slot it came from:
\begin{equation}
  (G^{\mathsf T}\vc)[n] \;=\; \sum_{(e,\ell)\,:\,\iota(e,\ell)=n} \vc[e,\ell].
  \label{eq:appk-scatter}
\end{equation}
This is exactly a \emph{scatter-add}: indexed writes that accumulate where indices
collide. It is the only stage in the pipeline that reduces---every other stage is a
map or a per-block contraction---and on a parallel machine it is the one stage that
needs an atomic or a coloring to make the colliding writes safe. The summation is
finite-element \emph{assembly} (\cref{ch:assembly}), here applied to a vector
rather than stamped into a matrix.

The composition $G^{\mathsf T}G$ is worth isolating, because it has a clean
meaning. Gather then immediately scatter-add: each global dof is copied out to
every element that touches it, then those copies are summed straight back with no
intervening computation. The result is the global dof scaled by how many elements
share it:
\begin{equation}
  (G^{\mathsf T}G\,\vv)[n] \;=\; m_n\,\vv[n],
  \qquad m_n = \#\{(e,\ell): \iota(e,\ell)=n\},
  \label{eq:appk-mult}
\end{equation}
the \emph{dof-multiplicity diagonal} $\mat{M}_{\mathrm{mult}}=\operatorname{diag}(m_n)$.
An interior vertex of a hex mesh has $m_n=8$; a face dof has $m_n=2$; a dof on the
domain boundary touched by one element has $m_n=1$. The multiplicity diagonal is
cheap, build-time, and useful: it is exactly the operator that converts a
\emph{summed} (assembled) quantity into an \emph{averaged} one, and it appears
whenever the pipeline must reconcile the redundant element-local representation
with the unique global one---for instance in building a diagonal for a smoother
(\cref{ch:smoothers}) or projecting between the two representations.

\begin{keyidea}
The gather $G$ and scatter $G^{\mathsf T}$ are a transpose pair built from one
index table $\iota(e,\ell)$: $G$ copies global dofs out to elements (duplicating
shared dofs), $G^{\mathsf T}$ sums element contributions back (assembling shared
dofs). All inter-element coupling in the whole operator lives in these two ends;
everything between them is block-diagonal in $E$. Their composition
$G^{\mathsf T}G=\operatorname{diag}(m_n)$ is the dof-multiplicity diagonal that
reconciles the duplicated element-local view with the unique global one.
\end{keyidea}

\section{Stage 2: \texorpdfstring{$B_{\diffop}$}{B}, basis evaluation}
\label{sec:appk-basis}

With the field restricted to elements, the weak form needs it---and its
derivative---sampled at the quadrature points. That is what $B_{\diffop}$ delivers:
it contracts each element's $L$ local coefficients against the tabulated reference
basis to produce $C$ components at each of $P$ quadrature points.

\subsection{The tabulated basis tensor}

The reference shape functions of \cref{ch:refelement} are sampled \emph{once} at
the quadrature nodes and stored. For the value (interpolation) mode this is a dense
matrix
\[
  \hat B \in \Real^{P\times L},
  \qquad \hat B_{q\ell} = \hat\phi_\ell(\hat\vx_q),
\]
the $\ell$-th reference shape function evaluated at the $q$-th quadrature point. For
a derivative mode the table is the reference \emph{gradient}
$\hat G\in\Real^{P\times L\times d}$ with $\hat G_{q\ell k}=\partial_k\hat\phi_\ell
(\hat\vx_q)$. Both are pure build-time data: tabulated once, shared by every element
of the mesh, independent of the input vector. The value-mode evaluation, per
element $e$, is the contraction over the local-dof axis $L$:
\begin{equation}
  (B\,\vc)[e,q,c] \;=\; \sum_{\ell=0}^{L-1} \hat B_{q\ell}\;\vc[e,\ell],
  \qquad C=1\ \text{for a scalar value},
  \label{eq:appk-basis}
\end{equation}
producing \code{[(E, P, C)]} from \code{[(E, L)]}. The same small $\hat B$ is
applied independently inside each of the $E$ elements---a batched matrix-vector
product, block-diagonal in $E$.

\subsection{The mode is selected by \texorpdfstring{$\diffop$}{D}}

The differential operator $\diffop$ of the weak-form term selects which tabulated
table $B_{\diffop}$ uses, and with it the meaning and extent of the component axis
$C$. \Cref{tab:appk-modes} lists the four de Rham modes (\cref{ch:refelement}).

\begin{table}[t]
\centering
\caption[Basis-evaluation modes and the component axis]{The four basis-evaluation
modes selected by the term's differential operator $\diffop$, with the
component-axis extent $C$ each produces. The interpolation mode evaluates basis
\emph{values} ($C=1$ scalar); the gradient, curl, and divergence modes evaluate
derivatives, with $C$ set by the de Rham slot. The same five-stage skeleton serves
all of them; only the tabulated table and the geometry-factor shape change.}
\label{tab:appk-modes}
\small
\setlength{\tabcolsep}{7pt}
\renewcommand{\arraystretch}{1.25}
\begin{tabular}{@{}llll@{}}
\toprule
$\diffop$ & mode & table & components $C$ \\
\midrule
$\mathrm{Id}$ & interp (values) & $\hat B$, $P\times L$ & $1$ (scalar) \\
$\Grad$ & grad & $\hat G$, $P\times L\times d$ & $d$ \\
$\Curl$ & curl & $\hat G$, curl-combined & $d$ (3-D), $1$ (2-D) \\
$\Div$ & div & $\hat G$, trace-combined & $1$ \\
\bottomrule
\end{tabular}
\end{table}

For a gradient term the contraction carries the extra $k$ axis,
\[
  (B_{\Grad}\vc)[e,q,c] \;=\; \sum_{\ell} \hat G_{q\ell c}\;\vc[e,\ell],
  \qquad c\in\{0,\dots,d-1\},
\]
so $C=d$: the reference gradient of the field at each quadrature point. (The
\emph{physical} gradient---the one the weak form actually integrates---needs the
inverse-transpose Jacobian $J^{-\mathsf T}$ applied to this reference gradient; that
multiplication is folded into the geometry factor and applied in stage 3, not here,
which keeps $B$ a pure reference-table contraction. We return to this in
\cref{sec:appk-geom}.) The curl and divergence modes combine the reference
derivatives the way the differential operator prescribes; the structural shape is
identical, only the table differs.

\subsection{The transpose $B_{\diffop}^{\mathsf T}$ integrates}

Stage 4 is the exact transpose. Where $B$ contracted the local-dof axis $L$ against
the table to produce quadrature-point values, $B_{\diffop}^{\mathsf T}$ contracts
the quadrature-point axis $P$ against the \emph{same} table to produce local-dof
contributions:
\begin{equation}
  (B^{\mathsf T}\vu)[e,\ell] \;=\; \sum_{q=0}^{P-1} \hat B_{q\ell}\;\vu[e,q,c],
  \label{eq:appk-basist}
\end{equation}
mapping \code{[(E, P, C)]} back to \code{[(E, L)]}. This is the discrete integral
$\int(\cdot)\,\phi_\ell$ against each test basis function: the quadrature sum
$\sum_q$ of \cref{ch:refelement} \emph{is} this transpose contraction, with the
quadrature weight already folded into the weighted values $\vu$ that reach it from
stage 3. The same tabulated $\hat B$ serves both directions---one of the reasons
the build-time stratum is so small.

\section{Stage 4 in passing, Stage 3 in full: the diagonal weight \texorpdfstring{$D$}{D}}
\label{sec:appk-diagonal}

We have already described stage 4 as the transpose of stage 2. The center stage
$D$---logically stage 3, the weight---is the simplest object in the pipeline and
the one that carries all the physics, so we give it its own treatment.

After basis evaluation the field is a tensor of values at quadrature points,
\code{[(E, P, C)]}. The weight stage multiplies each point's value, independently,
by a single precomputed factor read from \code{geom\_data}:
\begin{equation}
  (D\,\vu)[e,p,c'] \;=\; \sum_{c} \texttt{geom\_data}[e,p]\,{}_{c'c}\;\vu[e,p,c],
  \label{eq:appk-diagonal}
\end{equation}
where for a scalar (mass) term the factor is a single number and \cref{eq:appk-diagonal}
is a plain scaling; for a vector (grad--grad) term the factor is a small $d\times d$
block---the metric---and the sum over $c$ is a tiny per-point matrix-vector product
that may change the component count to $C'$. \Cref{fig:appk-diagonal} draws the
structure: the output at index $(e,p)$ depends only on the input at $(e,p)$ and on
\code{geom\_data}$[e,p]$, never on any other point.

\begin{figure}[t]
\centering
\begin{tikzpicture}[scale=1.0]
  \foreach \i in {0,...,4}{
    \foreach \j in {0,...,2}{
      \node[draw=simGreen, fill=simBgGreen, rounded corners=1pt,
            minimum size=7mm, font=\scriptsize] (c\i\j) at (\i*1.05, \j*0.9) {$u_{\i\j}$};
    }
  }
  \foreach \i in {0,...,4}{
    \foreach \j in {0,...,2}{
      \node[font=\scriptsize, simGold] at (\i*1.05+0.35, \j*0.9+0.35) {$\odot g$};
    }
  }
  \node[font=\scriptsize\itshape, simRust] at (-0.95, 0.9) {$(e,p)$};
  \node[font=\scriptsize, simInk] at (2.1,-0.85) {$P$ points $\times$ $E$ elements (flattened)};
  \draw[decorate, decoration={brace, amplitude=5pt}, simBlue]
    (-0.5, 2.2) -- (4.7, 2.2)
    node[midway, above=5pt, font=\footnotesize, simBlue]
    {every cell independent --- no coupling, any order, fully parallel};
\end{tikzpicture}
\caption[The weight stage as an embarrassingly parallel diagonal]{The weight stage
$D$ over the flattened $(e,p)$ index. Each cell is one quadrature point's evaluated
value $u_{ep}$, multiplied ($\odot$) by its own precomputed geometry block $g_{ep}$
from \code{geom\_data}. No cell depends on any other: $D$ is \emph{diagonal} in the
$(E,P)$ index---a per-point lift, mapped across the whole point set at once. There
is no reduction, no recurrence, no order-dependence; the output at $(e,p)$ is a
function of the input at $(e,p)$ alone. This is the embarrassingly parallel core of
the operator and the form a tensor accelerator consumes most directly.}
\label{fig:appk-diagonal}
\end{figure}

This is why we call $D$ the \emph{diagonal} of the operator. The basis stages
couple the $L$ local dofs to the $P$ quadrature points (a dense small contraction);
the gather and scatter couple elements through shared dofs. But $D$, sandwiched
between them, couples nothing: it is a diagonal scaling in the quadrature-point
basis, block-diagonal not merely in $E$ but in the full $(E,P)$ index. Everything
that distinguishes a mass term from a stiffness term, this material from that one,
a curved element from a straight one, is packed into the contents of
\code{geom\_data}---and that packing is the build-time pass we turn to now.

\section{Building the geometry factor: \texttt{geom\_factor\_build}}
\label{sec:appk-geom}

The tensor \code{geom\_data} of shape \code{[(E, P, G)]} is the one piece of the
build-time stratum that costs real arithmetic. It is produced by a setup pass,
\code{geom\_factor\_build}, whose signature is
\begin{lstlisting}[style=l4style]
-- the geometry-factor build pass: a pure setup-stratum function
geom_factor_build :: MeshNodes -> QuadWeights -> Tensor[(E, P, G)]
  -- per (e, p):  geom_data[e,p] = metric( J(mesh_nodes)[e,p] ) * w[e,p] * Q[e,p]
  --   E = elements;  P = quad points per element;  G = 2 + d*d metric components
\end{lstlisting}
At each quadrature point of each element it forms the Jacobian, combines it into
the metric the weak-form term needs, scales by the quadrature weight, and folds in
the material coefficient.

\subsection{From the Jacobian to the metric}

Recall from \cref{ch:refelement} that the geometric map $F_K$ carries the reference
element onto the physical element, and its derivative is the Jacobian $J=\partial
F_K/\partial\hat\vx\in\Real^{d\times d}$, whose columns are the physical edge
vectors. Two objects derived from $J$ are all the geometry a term needs:
\begin{equation}
\begin{aligned}
  \text{mass term ($\diffop=\mathrm{Id}$):}\quad & g = \abs{\det J}\,; \\
  \text{grad--grad ($\diffop=\Grad$):}\quad & \mat{G}_{\mathrm{metric}} = J^{-1}J^{-\mathsf T}\,\abs{\det J}.
\end{aligned}
  \label{eq:appk-metric}
\end{equation}
The mass factor is a scalar; the grad--grad metric is a symmetric $d\times d$
matrix. The metric arises by pulling both gradients in $\Grad\phi_i\cdot\Grad\phi_j$
back through $J^{-\mathsf T}$ and collecting the volume factor (\cref{ch:refelement}
derives it). Storing $J^{-1}J^{-\mathsf T}|\det J|$ is exactly the device that lets
the run-time $D$ apply the physical-gradient transformation as a single per-point
matrix multiply, with no Jacobian inversion at apply time.

\subsection{Folding in the weight and the coefficient}

The build pass does not stop at the bare metric. It pre-multiplies two more
factors, so that the run-time stage $D$ is a single multiply with nothing left to
combine:
\begin{equation}
  \texttt{geom\_data}[e,p] \;=\;
    \underbrace{\bigl(\text{metric from } J[e,p]\bigr)}_{\abs{\det J}\ \text{or}\ J^{-1}J^{-\mathsf T}\abs{\det J}}
    \;\cdot\;
    \underbrace{w_p}_{\text{quad weight}}
    \;\cdot\;
    \underbrace{Q[e,p]}_{\text{material coeff}}.
  \label{eq:appk-geomdata}
\end{equation}
The quadrature weight $w_p$ (\cref{ch:refelement}) rides inside \code{geom\_data}
rather than appearing as a separate stage---the weighted sum
$\sum_q w_q(\cdot)$ of the discrete integral is split across $D$ (the $w_q$ factor)
and $B^{\mathsf T}$ (the $\sum_q$ contraction). The material coefficient $Q$---the
$\varepsilon$, $\mu^{-1}$, or conductivity of the term, looked up by the element's
material-region attribute---is folded in the same way. The result is one number
(mass) or one small matrix (grad--grad) per quadrature point, $G=2+d^2$ stored
components, computed once and read by every apply.

\begin{remark}
On a \emph{straight-sided} (affine) element the Jacobian $J$ is constant over the
element, so the metric is constant in $p$ and \code{geom\_data}$[e,\cdot]$ repeats
the same block at every quadrature point---a degenerate case the build pass may
exploit by storing one block per element instead of $P$. On a \emph{curved}
(isoparametric high-order) element $J$ varies point to point and the full
\code{[(E, P, G)]} storage is needed. The pipeline is identical either way; only
the storage compresses.
\end{remark}

\section{Build-time versus run-time stratification}
\label{sec:appk-strata}

The single most consequential structural fact about the operator is that its data
divides into two strata by \emph{lifetime}. \Cref{ch:strata} treats this division
across the whole simulator; here we make it concrete for the kernel.

\begin{description}[leftmargin=1.4em, style=nextline]
\item[Build-time (setup) stratum.] Data fixed once per
$(\text{mesh},\text{FE order},\text{quadrature rule})$ and reused across every
application of the operator: the restriction index map $\iota(e,\ell)$, the
tabulated basis tables $\hat B$ and $\hat G$, and---the one with real
arithmetic---\code{geom\_data}. Their axis extents $E$, $L$, $P$, $G$ are all
build-time.
\item[Run-time (apply) stratum.] Data produced afresh on every action
$\vv\mapsto\mat{A}\vv$: the gathered \code{[(E, L)]} block, the evaluated and
weighted \code{[(E, P, C)]} values, the integrated contributions. These are
transient, born inside one apply and discarded when it returns. The only
run-time-\emph{shaped} axis is $C$, and even it is fixed once $\diffop$ is chosen;
what truly varies per apply is the field \emph{value} at each $(e,p,c)$ index.
\end{description}

\Cref{fig:appk-strata} draws the split. A setup pass runs the build \emph{once} and
produces the three tables; thereafter every operator application---and a Krylov
solve performs hundreds---reads them without rebuilding. The expensive geometry
work is amortized across the entire solve. Only when the mesh changes---adaptive
refinement (\cref{ch:amr}) being the usual cause---is the setup pass rerun, a
setup-stratum invalidation, not a per-apply cost.

\begin{figure}[t]
\centering
\begin{tikzpicture}[node distance=7mm and 12mm]
  % --- build-time row (gold) ---
  \node[data, fill=simBgGold, draw=simGold, text width=20mm] (mesh)
    {mesh nodes\\$+$ quad rule};
  \node[stage, fill=simBgGold, draw=simGold, right=of mesh, text width=27mm]
    (build) {\textbf{setup pass}\\\footnotesize $J,\;\abs{\det J},\;Q,\;w$\\(\code{geom\_factor\_build})};
  \node[data, fill=simBgGold, draw=simGold, right=of build, text width=23mm]
    (gd) {\code{geom\_data}\\\code{[(E, P, G)]}\\$+$ basis $+$ restr};
  \draw[flow, simGold] (mesh) -- (build);
  \draw[flow, simGold] (build) -- (gd);
  \node[font=\scriptsize\bfseries, simGold, left=1mm of mesh, rotate=90, anchor=south]
    {build-time};
  \node[font=\scriptsize\itshape, simGray, below=1mm of build]
    {runs \emph{once} per (mesh, order)};

  % --- run-time row (blue): many applies reading the tables ---
  \node[stage, below=18mm of build, text width=30mm] (apply1)
    {apply $\mat{A}\vv_1$\\\footnotesize $G^{\mathsf T}B^{\mathsf T}\,D\,B\,G$};
  \node[stage, right=5mm of apply1, text width=13mm] (apply2)
    {apply\\$\mat{A}\vv_2$};
  \node[font=\large, right=2mm of apply2] (dots) {$\cdots$};
  \node[stage, right=2mm of dots, text width=13mm] (applyk)
    {apply\\$\mat{A}\vv_k$};
  \node[font=\scriptsize\bfseries, simBlue, left=1mm of apply1, rotate=90, anchor=south]
    {run-time};
  \node[font=\scriptsize\itshape, simGray, below=1mm of apply1, xshift=14mm]
    {runs \emph{every} Krylov iteration --- reads the tables, never rebuilds them};
  \draw[flow, simGold] (gd.south) -- ++(0,-0.4) -| (apply1.north);
  \draw[flow, simGold] (gd.south) -- ++(0,-0.4) -| (apply2.north);
  \draw[flow, simGold] (gd.south) -- ++(0,-0.4) -| (applyk.north);
\end{tikzpicture}
\caption[Build-time versus run-time stratification]{The geometry factor and the two
index/basis tables are computed \emph{once}. A setup pass (gold, top) consumes the
mesh nodes and the quadrature rule and produces \code{geom\_data} together with the
restriction map and tabulated basis. Every subsequent operator application (blue,
bottom)---and an iterative solve performs hundreds---simply \emph{reads} these
build-time tables as it runs the five stages. The expensive geometry work is
amortized across the whole solve; only a mesh change (e.g.\ adaptive refinement,
\cref{ch:amr}) reruns the setup pass. This setup/apply split is the kernel's
instance of the stratification that organizes the whole simulator
(\cref{ch:strata}).}
\label{fig:appk-strata}
\end{figure}

\begin{keyidea}
The kernel carries two strata. The \emph{build-time} stratum---restriction map,
tabulated basis $\hat B,\hat G$, and the geometry factor \code{geom\_data}---is
computed once per $(\text{mesh},\text{order})$ and reused. The \emph{run-time}
stratum---the gathered, evaluated, weighted, integrated tensors---is rebuilt on
every application. Folding the Jacobian metric, quadrature weight, and material
coefficient into one precomputed \code{geom\_data} is precisely what reduces the
run-time weight stage $D$ to a single pointwise multiply.
\end{keyidea}

\section{Sum-factorization: the basis stage, reordered}
\label{sec:appk-sumfac}

The one stage with non-trivial arithmetic per apply is the basis evaluation $B$ and
its transpose. Taken naively, $B$ is a dense $P\times L$ contraction per element,
and at high order that is the dominant cost. On tensor-product elements
\emph{sum-factorization} cuts it dramatically by exploiting the product structure
of the basis---and it does so as a pure contraction reordering, computing
bit-for-bit the same result.

\subsection{The product structure}

On a tensor-product ($Q_p$) element the local dofs and the quadrature points each
lie on a $d$-dimensional grid of side $p+1$, so $L=P=(p+1)^d$, and the
$d$-dimensional basis factors into a product of one-dimensional bases:
\[
  \hat\phi_{i_1\cdots i_d}(\hat\xi_1,\dots,\hat\xi_d)
  \;=\; \hat\phi_{i_1}(\hat\xi_1)\,\cdots\,\hat\phi_{i_d}(\hat\xi_d).
\]
The naive evaluation ignores this and contracts against the full dense table:
\begin{equation}
  \text{naive cost} \;=\; P\cdot L \;=\; (p+1)^d\cdot(p+1)^d \;=\; (p+1)^{2d}
  \quad\text{multiply-adds per element.}
  \label{eq:appk-naive}
\end{equation}
Sum-factorization instead applies the \emph{one}-dimensional basis along each axis
in turn. Contract axis $1$ (a $(p+1)\times(p+1)$ matrix applied along one grid
direction, the other $d-1$ directions carried along as batch), then axis $2$, and
so on. Each of the $d$ sweeps costs $(p+1)$---the 1-D contraction---times
$(p+1)^d$---the rest of the grid carried along:
\begin{equation}
  \text{sum-factored cost} \;=\; d\cdot(p+1)\cdot(p+1)^d \;=\; d\,(p+1)^{d+1}.
  \label{eq:appk-factored}
\end{equation}
In the operator-count notation of \cref{ch:matrixfree} this is the move from
$O(P\cdot L)$ to $O(d\cdot P^{1/d}\cdot L)$, since one grid axis has length
$P^{1/d}=(p+1)$. \Cref{fig:appk-sumfac} draws the reordering for $d=3$ as a chain
of three batched 1-D contractions.

\begin{figure}[t]
\centering
\begin{tikzpicture}[node distance=5mm and 9mm]
  % naive: one dense contraction
  \node[data, text width=16mm] (din) {dofs\\$(p{+}1)^3$};
  \node[opbox, fill=simBgRust, draw=simRust, right=of din, text width=26mm]
    (naive) {dense $P\times L$\\\footnotesize $(p{+}1)^6$ ops};
  \node[data, right=of naive, text width=16mm] (dpts) {pts\\$(p{+}1)^3$};
  \draw[flow] (din) -- (naive);
  \draw[flow] (naive) -- (dpts);
  \node[font=\scriptsize\bfseries, simRust, left=1mm of din, rotate=90, anchor=south]
    {naive};

  % factored: three 1-D sweeps (d=3)
  \node[data, below=15mm of din, text width=16mm] (fin) {dofs\\$(p{+}1)^3$};
  \node[opbox, right=of fin, text width=17mm] (s1)
    {sweep $\xi_3$\\\footnotesize $(p{+}1)^4$};
  \node[opbox, right=of s1, text width=17mm] (s2)
    {sweep $\xi_2$\\\footnotesize $(p{+}1)^4$};
  \node[opbox, right=of s2, text width=17mm] (s3)
    {sweep $\xi_1$\\\footnotesize $(p{+}1)^4$};
  \node[data, right=of s3, text width=16mm] (fpts) {pts\\$(p{+}1)^3$};
  \draw[flow] (fin) -- (s1);
  \draw[flow] (s1) -- node[above, font=\scriptsize, simGray] {partial} (s2);
  \draw[flow] (s2) -- node[above, font=\scriptsize, simGray] {partial} (s3);
  \draw[flow] (s3) -- (fpts);
  \node[font=\scriptsize\bfseries, simBlue, left=1mm of fin, rotate=90, anchor=south]
    {factored};
\end{tikzpicture}
\caption[Sum-factorization as a contraction reordering ($d=3$)]{Basis evaluation on
a 3-D tensor-product element, two ways. \emph{Top (rust):} the naive contraction
treats the basis as one dense $P\times L$ matrix, costing $(p+1)^6$ multiply-adds
per element. \emph{Bottom (blue):} sum-factorization applies the one-dimensional
basis along each of the three axes in turn---a sweep in $\xi_3$, then $\xi_2$, then
$\xi_1$, each carrying the other two grid directions as batch and costing
$(p+1)^4$---for $3(p+1)^4$ total. The two produce \emph{identical} values; only the
contraction order differs. The saving is the move from $O(P\cdot L)$ to
$O(d\,P^{1/d}L)$, and the exponent gap ($2d$ versus $d+1$) makes it decisive at
high order.}
\label{fig:appk-sumfac}
\end{figure}

\begin{table}[t]
\centering
\caption[Sum-factorization operation counts, $d=3$]{Multiply-adds per element for
the basis-evaluation stage in three dimensions, naive $(p+1)^{2d}=(p+1)^6$ versus
sum-factored $d\,(p+1)^{d+1}=3(p+1)^4$. The speedup is unbounded in $p$: the
exponent drops from $2d=6$ to $d+1=4$, so the ratio grows like $(p+1)^{d-1}/d$.}
\label{tab:appk-sumfac}
\small
\setlength{\tabcolsep}{9pt}
\renewcommand{\arraystretch}{1.2}
\begin{tabular}{@{}cccc@{}}
\toprule
order $p$ & naive $(p{+}1)^6$ & factored $3(p{+}1)^4$ & speedup \\
\midrule
1 & 64               & 48                & $1.3\times$ \\
2 & 729              & 243               & $3.0\times$ \\
4 & \num{15625}      & \num{1875}        & $8.3\times$ \\
8 & \num{531441}     & \num{19683}       & $27\times$ \\
\bottomrule
\end{tabular}
\end{table}

\Cref{tab:appk-sumfac} tabulates the 3-D savings. At $p=8$ it is the difference
between $9^6=\num{531441}$ and $3\cdot 9^4=\num{19683}$ operations per element---a
factor of $27$, widening with every increment of $p$. Without sum-factorization the
basis stage would reintroduce the very $O(p^{2d})$ cost the matrix-free form set out
to escape; with it, high-order matrix-free is not merely affordable but the fastest
form available.

\begin{pitfall}
Sum-factorization is a \emph{transparent} performance trick (\cref{ch:bigidea}):
same result, faster, with no change to the answer, the conditioning, or the
rounding behavior that the algorithm depends on. It is recorded as a one-line note
on the basis stage, never as a separate algebraic claim. Do not mistake it for a
\emph{load-bearing} numerical trick---it reorders the sums of a single linear map
and nothing more. The reason to know it is that it is \emph{why} the high-order
kernel is fast, not a change to what the kernel computes.
\end{pitfall}

\section{The libCEED boundary: API versus implementation}
\label{sec:appk-extern}

We have now specified all five stages as tensor contractions. One honest
qualification remains, and it is the subject of this section: in the real Palace
build the innermost element-quadrature kernel is not hand-written from these
contractions but handed to a specialized accelerator library, libCEED, which owns
the device-level details---which loops to fuse, how to tile the quadrature points,
how to map them onto warps or wavefronts. The book draws this as a clean boundary
with two sides, and the distinction between them is load-bearing.

\subsection{The two surfaces}

The boundary has a \emph{kernel API} and a \emph{kernel implementation}, and the
spec keeps both, linked for review (\cref{ch:synthlib}).

\paragraph{The kernel API---what we call.} The API surface is the opaque library
contract: the shape signature and the semantics of the innermost kernel, with
\emph{no} claim about how it is realized. In the calculus it is written with the
\code{\#extern} marker---a node whose contract we state precisely but whose body we
delegate:
\begin{lstlisting}[style=l4style]
-- the kernel API: the contract is ours, the device body is libCEED's
quad_kernel :: GeomData -> Tensor[(E, P, C)] -> Tensor[(E, P, C')]
#extern quad_kernel   -- shape + semantics specified; body delegated
\end{lstlisting}
This is the violet, dashed \code{\#extern} node at the center of
\cref{fig:appk-pipeline}. Its contract is exactly the pointwise diagonal of
\cref{sec:appk-diagonal}: a shape-preserving (up to $C\to C'$) elementwise multiply
against \code{geom\_data}. The library owns the machine; we own the contract.

\paragraph{The kernel implementation---the contraction chain.} The implementation
surface is the constructive realization in our own tensor vocabulary: the five-stage
chain $G^{\mathsf T}B_{\diffop}^{\mathsf T}D B_{\diffop}G$ of this entire appendix,
built from the firm substrate operators \code{element\_restrict}, \code{basis\_apply},
\code{quad\_point\_contract}, and \code{geom\_factor\_build}. It makes a positive
claim and carries a normal rank: it \emph{is} the kernel, expressed from our
primitives.

\begin{lstlisting}[style=l4style]
-- the kernel implementation: the contraction chain that realizes quad_kernel,
-- written from our own firm substrate operators.  apply A v = scatter . Bt . D . B . gather
apply_kernel :: Operator -> Tensor[(N: ..)] -> Tensor[(N: ..)]
apply_kernel op v =
  scatter   op.restr           $   -- G^T  : [(E, L)]      -> [(N: ..)]   element_restrict^T
  integrate op.mode op.basis   $   -- B_D^T: [(E, P, C')]  -> [(E, L)]    basis_apply^T
  weight    op.geom            $   -- D    : [(E, P, C)]   -> [(E, P, C')] quad_point_contract  (= #extern quad_kernel)
  evaluate  op.mode op.basis   $   -- B_D  : [(E, L)]      -> [(E, P, C)] basis_apply
  gather    op.restr v             -- G    : [(N: ..)]     -> [(E, L)]    element_restrict
  where
    -- op.geom is the build-time geom_data, computed once:
    -- op.geom = geom_factor_build mesh_nodes quad_weights
\end{lstlisting}

\subsection{The correspondence}

The two surfaces are tied by a typed \emph{\code{realizes-kernel-api}}
correspondence: an edge from the implementation node to the API node asserting
``this contraction chain realizes that opaque contract.'' \Cref{fig:appk-extern}
draws it. The edge is a \emph{reference}, not a dependency: the implementation does
not build on the opaque API---it is an \emph{alternative} to it---so the edge
carries no rank or liveness constraint. A reviewer reads both sides and checks they
match: the opaque \code{\#extern} contract against the from-our-primitives chain.
That review is the discipline the boundary buys; the lowering-verifier audits the
correspondence.

\begin{figure}[t]
\centering
\resizebox{\linewidth}{!}{%
\begin{tikzpicture}[node distance=9mm and 22mm]
  % API node (extern, violet)
  \node[extern, text width=42mm] (api)
    {\textbf{kernel API}\;\code{\#extern}\\[2pt]
     \footnotesize \code{quad\_kernel ::}\\
     \footnotesize \code{GeomData -> [(E,P,C)] -> [(E,P,C')]}\\[2pt]
     \footnotesize\itshape opaque libCEED contract\\(shape $+$ semantics; body delegated)};
  % impl node (stage, blue)
  \node[stage, below=of api, text width=58mm] (impl)
    {\textbf{kernel implementation}\\[2pt]
     \footnotesize $G^{\mathsf T}\,B_{\diffop}^{\mathsf T}\,D\,B_{\diffop}\,G$\\[2pt]
     \footnotesize from firm substrate:\\
     \footnotesize \code{element\_restrict}, \code{basis\_apply},\\
     \footnotesize \code{quad\_point\_contract}, \code{geom\_factor\_build}};
  % the realizes-kernel-api reference edge (free, navigational)
  \draw[refedge] (impl.north) -- node[right, font=\scriptsize, simGray, align=left]
     {\code{realizes-}\\\code{kernel-api}\\(reference:\\free, reviewed)} (api.south);
  % the substrate, to the side
  \node[opbox, right=34mm of impl, text width=30mm] (sub)
    {\footnotesize firm \code{[(E,L)]}/\allowbreak\code{[(E,P,C)]}/\allowbreak\code{[(E,P,G)]}\\substrate ops};
  \draw[depedge] (impl.east) -- node[above, font=\scriptsize, simBlue] {composes} (sub.west);
\end{tikzpicture}%
}
\caption[The kernel API and the kernel implementation]{The two surfaces of the
libCEED boundary, kept distinct and linked for review (\cref{ch:synthlib}).
\emph{Top (violet, \code{\#extern}):} the kernel \emph{API}---the opaque libCEED
contract, stating the shape signature and pointwise semantics of the innermost
kernel with no claim about its realization. \emph{Bottom (blue):} the kernel
\emph{implementation}---the five-stage contraction chain of this appendix, built by
composing the firm \code{[(E,L)]}/\code{[(E,P,C)]}/\code{[(E,P,G)]} substrate
operators. The \code{realizes-kernel-api} edge (dotted) is a \emph{reference}, not a
dependency: the implementation is an alternative to the opaque API, not a consumer
of it, so the edge carries no rank or liveness---it is a correspondence a reviewer
checks. The same \code{\#extern} pattern recurs at the SLEPc eigensolve
(\cref{ch:eigenvalues}) and the triangular solves inside a smoother
(\cref{ch:smoothers}).}
\label{fig:appk-extern}
\end{figure}

\begin{keyidea}
The matrix-free operator is a chain of element-local tensor contractions over named
axes---the GPU-native form, and the kernel \emph{implementation} behind the libCEED
boundary. The boundary has two surfaces: the kernel \emph{API}, an opaque
\code{\#extern} contract stating shape and semantics; and the kernel
\emph{implementation}, the chain $G^{\mathsf T}B_{\diffop}^{\mathsf T}D B_{\diffop}G$
built from firm substrate operators. A \code{realizes-kernel-api} reference edge
links them for review. We own the mathematics; the library owns the machine; and
the two meet at a narrow, well-typed, reviewable interface.
\end{keyidea}

\section{Why this is the accelerator target}
\label{sec:appk-gpu}

Step back from the stages and look at what the whole operator turned out to be. The
finite-element operator $\mat{A}$---the discretization of Maxwell's weak form, the
object every solver in the book applies---is a composition of tensor contractions
over named axes. Gather is an indexed copy. Evaluate and integrate are dense
contractions of an \code{[(E, L)]} tensor against a tabulated basis, yielding
\code{[(E, P, C)]}, and on tensor-product elements they factor further into batched
1-D contractions. Weight is an elementwise multiply of \code{[(E, P, C)]} against
\code{[(E, P, G)]}. Scatter is an indexed sum. There is no sparse-matrix data
structure anywhere---no row pointers, no column indices, no stored entries. There
are only tensors and contractions.

This is precisely the shape a tensor accelerator is built to run. Every stage is a
pure function from tensors to tensors with no aliasing and no ordering hazard; the
element axis $E$ is a large batch dimension the hardware maps directly onto its
parallel units; the weight stage is a perfectly diagonal map; and the only stages
that couple data---gather and scatter---are pure index movement at the two ends,
not arithmetic in the hot loop. The assembled-matrix form, by contrast, is a
CPU-era data structure: irregular, pointer-chasing, bandwidth-bound, something a
tensor backend must be coaxed into. The matrix-free kernel is the form the
accelerator was built for, and it is exactly what the book's implementation view
(\cref{ch:synthlib}) expands the ``\,\lstinline[style=l4style]|apply A v|\,'' of
every solver into.

\summarysection
This appendix opened the five factors of $\mat{A}\vv=
G^{\mathsf T}B_{\diffop}^{\mathsf T}\,D\,B_{\diffop}\,G\,\vv$ down to their exact
tensor shapes. The gather $G$ is a pure indexed copy
$\texttt{[(N: \dots)]}\to\texttt{[(E, L)]}$ defined by the element-to-dof index map
$\iota(e,\ell)$; its transpose $G^{\mathsf T}$ is the scatter-add (assembly), the
operator's only reduction, and the composition $G^{\mathsf T}G=\operatorname{diag}
(m_n)$ is the dof-multiplicity diagonal. Basis evaluation $B_{\diffop}$ contracts
the local-dof axis $L$ against the tabulated reference basis to produce
\code{[(E, P, C)]}, with the mode---interp, grad, curl, div---and the component
count $C$ selected by the term's differential operator $\diffop$; $B^{\mathsf T}$
contracts the quadrature axis $P$ back to local dofs, carrying the discrete
integral. The center stage $D$ is the embarrassingly parallel diagonal: a pointwise
multiply against \code{geom\_data}, block-diagonal in the full $(E,P)$ index. That
geometry factor, of shape \code{[(E, P, G)]} with $G=2+d^2$ components, is built
once by \code{geom\_factor\_build} from the Jacobian metric
($\abs{\det J}$ for mass, $J^{-1}J^{-\mathsf T}\abs{\det J}$ for grad--grad), with
the quadrature weight and material coefficient pre-multiplied in. This is the
build-time stratum---computed once per $(\text{mesh},\text{order})$ and reused
across the hundreds of applies a solve performs---as against the transient run-time
tensors rebuilt on every action. On tensor-product elements sum-factorization
reorders the basis contraction from $O(P\cdot L)=O((p+1)^{2d})$ to
$O(d\,P^{1/d}L)=O(d(p+1)^{d+1})$---a transparent reordering, same answer, $27\times$
at $p=8$ in 3-D. The innermost kernel is sealed behind the libCEED boundary, which
the spec gives two surfaces: an opaque \code{\#extern} kernel \emph{API} stating
shape and semantics, and a kernel \emph{implementation}---the contraction chain of
this appendix from firm substrate operators---tied to it by a
\code{realizes-kernel-api} reference edge for review. Because every stage is a
contraction over named axes with no sparse-matrix structure anywhere, the
matrix-free kernel \emph{is} the GPU-native implementation view (\cref{ch:synthlib})
that the book has been building toward.

\furtherreading
The matrix-free, sum-factorized, tensor-contraction view of high-order
finite-element operators, and the kernel-API/implementation split realized through
libCEED's \code{\#extern} boundary, are the design philosophy of the CEED project;
its co-design survey~\citep{ceed2021} is the entry point to the libraries (libCEED,
MFEM) that realize the five-stage decomposition and the operation-count argument for
sum-factorization. For the finite-element foundations this appendix builds on---the
reference element, the tabulated bases and quadrature, the Jacobian metric, and the
element matrices the operator factors---Jin's text on the finite-element method in
electromagnetics~\citep{jin2014} gives the Maxwell-specific grounding, and the
mathematical theory of the spaces and their approximation properties is developed in
Brenner and Scott~\citep{brennerscott2008}. The build-time/run-time stratification
the geometry factor relies on is treated in its own right in \cref{ch:strata}; the
implementation view the contraction chain feeds is \cref{ch:synthlib}.
